\documentclass[12pt]{article}

% ============================================================
% Compile: pdflatex (mặc định của Overleaf) hoặc xelatex đều được.
% Cần file logo-bk.png đặt cùng thư mục với file .tex.
% Hình EDA (eda.py -> output/eda/) và hình A1.2 (analyze.py -> output/a1_2/).
% Bảng và số liệu A1.2 nằm trong generated/ (analyze.py sinh ra); mọi con số trong
% phần kết quả được chèn bằng \val{...}, nên báo cáo luôn khớp với log huấn luyện.
% ============================================================
\usepackage{iftex}
\ifPDFTeX
  \usepackage[utf8]{inputenc}
  \usepackage[T5]{fontenc}
  \usepackage{lmodern}
\else
  \usepackage{fontspec}   % XeLaTeX/LuaLaTeX: Latin Modern (Unicode) mặc định
\fi
\usepackage[vietnamese]{babel}

\usepackage[a4paper,margin=2.5cm,headheight=2cm,headsep=0.8cm]{geometry}
\usepackage{graphicx,booktabs,tabularx,xltabular,longtable,array,xcolor,enumitem,amssymb,multirow,caption,makecell}
\usepackage{fancyhdr}
\usepackage{lastpage}
\usepackage{titlesec}
\usepackage{float}
\usepackage{hyperref}
\hypersetup{hidelinks, pdftitle={A1 - CNN vs Transformer tren CIFAR-100}}
\usepackage{parskip}

\graphicspath{{./}{figures/}{../output/eda/}{../output/a1_2/}}
\captionsetup{labelfont=bf,skip=4pt}
\setlist{nosep,leftmargin=*}
\titleformat{\section}{\large\bfseries}{\thesection.}{0.4em}{}
\titleformat{\subsection}{\normalsize\bfseries}{\thesubsection.}{0.4em}{}
\titlespacing*{\section}{0pt}{10pt}{4pt}\titlespacing*{\subsection}{0pt}{6pt}{2pt}
\newcolumntype{L}{>{\raggedright\arraybackslash}X}
\newcolumntype{C}{>{\centering\arraybackslash}X}
\newcolumntype{P}[1]{>{\raggedright\arraybackslash}p{#1}}
\newcolumntype{R}[1]{>{\raggedleft\arraybackslash}p{#1}}
\renewcommand{\arraystretch}{1.15}
\setlength{\emergencystretch}{3em}
\newcommand{\code}[1]{\texttt{#1}}
\newcommand{\st}[1]{\textbf{#1}}
\newcommand{\todo}[1]{\textcolor{red}{\textbf{[#1]}}}
% Chèn hình nếu có file, nếu chưa upload thì hiện khung báo thiếu (để vẫn compile được)
\newcommand{\fig}[2]{\IfFileExists{#2}{\includegraphics[width=#1]{#2}}{%
  \IfFileExists{../output/eda/#2}{\includegraphics[width=#1]{../output/eda/#2}}{%
  \IfFileExists{../output/a1_2/#2}{\includegraphics[width=#1]{../output/a1_2/#2}}{%
  \fbox{\parbox{0.8\linewidth}{\centering\todo{Thiếu hình \detokenize{#2}}}}}}}}
% Số liệu A1.2 do analyze.py sinh ra (generated/macros.tex). \val{ResNetFullTest} in giá trị,
% hoặc ô đỏ [?] nếu chưa có kết quả, để báo cáo vẫn compile được khi chưa chạy xong.
\newif\ifresults
\IfFileExists{generated/macros.tex}{\resultstrue\newcommand{\ResNetScratchTest}{67,46}
\newcommand{\ResNetScratchFone}{67,50}
\newcommand{\ResNetScratchVal}{68,54}
\newcommand{\ResNetScratchTrainAcc}{98,59}
\newcommand{\ResNetScratchMin}{37,0}
\newcommand{\ResNetScratchSecEpoch}{74}
\newcommand{\ResNetScratchMem}{1,69}
\newcommand{\ResNetScratchBestEpoch}{27}
\newcommand{\ResNetHeadTest}{51,02}
\newcommand{\ResNetHeadFone}{49,24}
\newcommand{\ResNetHeadVal}{50,82}
\newcommand{\ResNetHeadTrainAcc}{51,17}
\newcommand{\ResNetHeadMin}{3,9}
\newcommand{\ResNetHeadSecEpoch}{23}
\newcommand{\ResNetHeadMem}{0,60}
\newcommand{\ResNetHeadBestEpoch}{10}
\newcommand{\ResNetPartialTest}{73,67}
\newcommand{\ResNetPartialFone}{73,42}
\newcommand{\ResNetPartialVal}{74,42}
\newcommand{\ResNetPartialTrainAcc}{76,03}
\newcommand{\ResNetPartialMin}{6,0}
\newcommand{\ResNetPartialSecEpoch}{36}
\newcommand{\ResNetPartialMem}{0,72}
\newcommand{\ResNetPartialBestEpoch}{9}
\newcommand{\ResNetFullTest}{74,71}
\newcommand{\ResNetFullFone}{74,48}
\newcommand{\ResNetFullVal}{74,92}
\newcommand{\ResNetFullTrainAcc}{77,04}
\newcommand{\ResNetFullMin}{12,3}
\newcommand{\ResNetFullSecEpoch}{74}
\newcommand{\ResNetFullMem}{1,69}
\newcommand{\ResNetFullBestEpoch}{9}
\newcommand{\DeiTScratchTest}{48,48}
\newcommand{\DeiTScratchFone}{48,01}
\newcommand{\DeiTScratchVal}{47,86}
\newcommand{\DeiTScratchTrainAcc}{96,49}
\newcommand{\DeiTScratchMin}{32,9}
\newcommand{\DeiTScratchSecEpoch}{66}
\newcommand{\DeiTScratchMem}{2,23}
\newcommand{\DeiTScratchBestEpoch}{20}
\newcommand{\DeiTHeadTest}{65,11}
\newcommand{\DeiTHeadFone}{64,72}
\newcommand{\DeiTHeadVal}{65,62}
\newcommand{\DeiTHeadTrainAcc}{65,16}
\newcommand{\DeiTHeadMin}{3,9}
\newcommand{\DeiTHeadSecEpoch}{23}
\newcommand{\DeiTHeadMem}{0,25}
\newcommand{\DeiTHeadBestEpoch}{10}
\newcommand{\DeiTPartialTest}{81,31}
\newcommand{\DeiTPartialFone}{81,30}
\newcommand{\DeiTPartialVal}{80,76}
\newcommand{\DeiTPartialTrainAcc}{93,40}
\newcommand{\DeiTPartialMin}{7,2}
\newcommand{\DeiTPartialSecEpoch}{43}
\newcommand{\DeiTPartialMem}{1,19}
\newcommand{\DeiTPartialBestEpoch}{7}
\newcommand{\DeiTFullTest}{85,02}
\newcommand{\DeiTFullFone}{85,04}
\newcommand{\DeiTFullVal}{84,94}
\newcommand{\DeiTFullTrainAcc}{97,18}
\newcommand{\DeiTFullMin}{11,0}
\newcommand{\DeiTFullSecEpoch}{66}
\newcommand{\DeiTFullMem}{2,23}
\newcommand{\DeiTFullBestEpoch}{10}
\newcommand{\BestResNetName}{\code{resnet18\_full}}
\newcommand{\BestResNetSameSuper}{44,4}
\newcommand{\BestResNetCoarseAcc}{85,93}
\newcommand{\BestResNetWorst}{girl (31\%), otter (40\%), seal (41\%), beaver (41\%), willow tree (44\%)}
\newcommand{\BestResNetTop}{motorcycle (96\%), wardrobe (94\%), bicycle (93\%), orange (93\%), television (93\%)}
\newcommand{\BestDeiTName}{\code{deit\_tiny\_full}}
\newcommand{\BestDeiTSameSuper}{48,5}
\newcommand{\BestDeiTCoarseAcc}{92,28}
\newcommand{\BestDeiTWorst}{girl (58\%), boy (58\%), man (66\%), maple tree (67\%), woman (69\%)}
\newcommand{\BestDeiTTop}{skunk (100\%), road (99\%), tractor (98\%), motorcycle (97\%), pickup truck (96\%)}
\newcommand{\AgreeBothRight}{70,59}
\newcommand{\AgreeOnlyCnn}{4,12}
\newcommand{\AgreeOnlyVit}{14,43}
\newcommand{\AgreeBothWrong}{10,86}
\newcommand{\OracleAcc}{89,14}
\newcommand{\RunsDone}{8}
\newcommand{\GpuName}{Tesla T4}
}{}
\newcommand{\val}[1]{\ifcsname #1\endcsname\csname #1\endcsname\else\todo{?}\fi}
% \gen{file}{số cột}: chèn các dòng bảng do analyze.py sinh, hoặc một dòng báo thiếu
% (dùng \@@input của TeX vì \input của LaTeX chèn thêm lệnh làm hỏng \bottomrule ngay sau bảng)
\makeatletter
\newcommand{\gen}[2]{\ifresults\@@input generated/#1 \else\multicolumn{#2}{c}{\todo{Chưa có kết quả}}\\\fi}
\makeatother

% ---------- Thông tin chung ----------
\newcommand{\TenTruong}{Đại học Quốc gia Thành phố Hồ Chí Minh\\Trường Đại học Bách khoa\\Khoa Khoa học và Kỹ thuật Máy tính}
\newcommand{\TenTruongHeader}{Ho Chi Minh University of Technology\\Faculty of Computer Science and Engineering}
\newcommand{\TenMonHoc}{Học sâu và ứng dụng trong thị giác máy tính}
\newcommand{\MaMonHoc}{CO5085}
\newcommand{\TenBaoCao}{Bài tập lớn A1 - Báo cáo cuối (A1.1 + A1.2)}
\newcommand{\TenDeTai}{CNN hay Transformer? So sánh ResNet-18 và DeiT-Tiny\\trên CIFAR-100 với nhiều chế độ huấn luyện}
\newcommand{\GiangVien}{Thầy Lê Thành Sách\\Thầy Nguyễn Quốc Minh}
\newcommand{\ThanhVien}{%
  Đàm Quang Phục - 2570483\\
  Hồ Bảo An - 2570164
}
\newcommand{\DiaDiemNgay}{Thành Phố Hồ Chí Minh, 11/2026}
\newcommand{\FooterText}{\TenMonHoc\ -- \MaMonHoc}

% ---------- Header / Footer ----------
\pagestyle{fancy}
\fancyhf{}
\renewcommand{\headrulewidth}{0.4pt}
\renewcommand{\footrulewidth}{0.4pt}
\fancyhead[L]{%
  \parbox[t][\headheight][t]{\headwidth}{%
    \vspace{0.2cm}
    \begin{minipage}[b]{1.1cm}
      \IfFileExists{logo-bk.png}{\includegraphics[width=1.1cm]{logo-bk.png}}{}
    \end{minipage}%
    \hspace{0.5em}%
    \begin{minipage}[b]{8cm}
      \small \TenTruongHeader
    \end{minipage}%
  }%
}
\fancyfoot[L]{\small \FooterText}
\fancyfoot[R]{\small Page \thepage/\pageref{LastPage}}

\fancypagestyle{plain}{%
  \fancyhf{}
  \renewcommand{\headrulewidth}{0pt}
  \renewcommand{\footrulewidth}{0pt}
}

\begin{document}

% ================= TRANG BÌA =================
\begin{titlepage}
\thispagestyle{plain}
\begin{center}
{\bfseries \TenTruong \par}
\vspace{2em}

\IfFileExists{logo-bk.png}{\includegraphics[width=3cm]{logo-bk.png}\\[1em]}{}

\vspace{5em}

{\LARGE \bfseries \TenMonHoc\ (\MaMonHoc) \par}
\vspace{0.3em}
\rule{0.85\textwidth}{0.5pt}

\vspace{1em}
{\Large \bfseries \TenBaoCao \par}
\vspace{0.6em}
{\large \TenDeTai \par}
\vspace{0.3em}
\rule{0.85\textwidth}{0.5pt}

\vspace{5em}

\noindent\begin{minipage}[t]{3.6cm}
\raggedleft
Giảng viên:\\[0.5em]
\phantom{x}\\[0.5em]
Thành viên nhóm:
\end{minipage}%
\hspace{0.6em}%
\begin{minipage}[t]{8cm}
\GiangVien\\[0.5em]
\ThanhVien
\end{minipage}

\vspace{2em}
{\footnotesize Mã nguồn: \url{https://github.com/dqpinfo/Deep_Learning_Computer_Vision_CO5085/tree/main/Assignment_1}\par}

\vfill
{\bfseries \DiaDiemNgay}
\end{center}
\end{titlepage}

% ================= MỤC LỤC =================
\tableofcontents
\newpage


% ================= MỤC LỤC =================
\tableofcontents
\newpage

% ======================================================================= status box
\noindent\fcolorbox{black!35}{black!3}{\parbox{\dimexpr\linewidth-2\fboxsep-2\fboxrule}{
\textbf{Tóm tắt.} Nhóm so sánh \st{ResNet-18} (CNN) và \st{DeiT-Tiny} (Transformer) trên \st{CIFAR-100}
(45.000 / 5.000 / 10.000 ảnh train / validation / test, split cố định seed 42) qua \st{8 thí nghiệm}: 2 kiến trúc
$\times$ 4 chế độ (học từ đầu, chỉ train head, đóng băng một phần, fine-tune toàn bộ). Mục 1--6 là phần A1.1
(dataset và protocol, hạn 26/10/2026); Mục 7--11 là phần A1.2 (kết quả và thảo luận, hạn 09/11/2026).
Kết quả chính: cấu hình tốt nhất là \code{deit\_tiny\_full} (test accuracy \val{DeiTFullTest}\%, 11 phút trên T4),
tiếp theo là \code{deit\_tiny\_partial} (\val{DeiTPartialTest}\%) và \code{resnet18\_full} (\val{ResNetFullTest}\%).
Pretrain quan trọng với ViT hơn hẳn với CNN: học từ đầu, DeiT-Tiny chỉ đạt \val{DeiTScratchTest}\% (kém ResNet-18
\val{ResNetScratchTest}\%), nhưng khi dùng trọng số pretrained, DeiT-Tiny vượt ResNet-18 ở cả ba mức đóng băng.
Đóng băng một phần gần bằng fine-tune toàn bộ với ResNet-18 (kém 1,0 điểm, nhanh gấp đôi) nhưng kém 3,7 điểm với
DeiT-Tiny; chỉ train head kém xa ở cả hai họ.
Mọi con số A1.2 được \code{analyze.py} sinh tự động từ log huấn luyện (\code{output/runs/*.json}).}}

% ======================================================================= 1
\section{Chủ đề và mục tiêu}
Bài A1 so sánh hai họ kiến trúc cho bài toán phân loại ảnh: mạng tích chập (CNN) và Vision Transformer (ViT)~\cite{dosovitskiy2021}.
Hai họ được đánh giá dưới nhiều chế độ huấn luyện: học từ đầu (\emph{from-scratch}) so với dùng trọng số
\emph{pretrained}, và ba mức đóng băng backbone trên bản pretrained (chỉ train head, đóng băng một phần,
fine-tune toàn bộ).

Báo cáo gồm hai phần. Phần A1.1 (Mục 1--6) chốt tập dữ liệu, cách chia split, DataLoader và augmentation, các mô
hình sử dụng và protocol thí nghiệm. Phần A1.2 (Mục 7--11) trình bày môi trường chạy, kết quả 8 thí nghiệm, phân tích
lỗi và thảo luận. Protocol ở phần A1.1 được giữ nguyên khi chạy A1.2; các thay đổi nhỏ về cách cài đặt (không đổi
protocol) được ghi ở Mục~\ref{sec:env}.

% ======================================================================= 2
\section{Tập dữ liệu: CIFAR-100}
\subsection{Lý do chọn}
\begin{itemize}
  \item \st{Đúng ràng buộc của đề:} tập ảnh công khai, có 100 lớp (yêu cầu $\geq 10$) và 50.000 ảnh huấn luyện;
        không trùng với các tập đã dùng ở E1 (MNIST / Fashion-MNIST / CIFAR-10).
  \item \st{Đủ khó để so sánh có ý nghĩa:} mỗi lớp chỉ có 500 ảnh train, ảnh nhỏ $32\times32$ và nhiều lớp rất giống
        nhau (ví dụ 5 loại cây: maple, oak, palm, pine, willow). Trong điều kiện này, ưu thế của pretrained và nhu cầu
        dữ liệu của Transformer bộc lộ rõ.
  \item \st{Vừa sức tính toán:} toàn bộ tập khoảng 140\,MB; cả 8 cấu hình có thể chạy trên một GPU phổ thông
        (Google Colab T4) trong thời gian của mốc A1.2.
  \item \st{Có nhãn hai cấp:} 100 lớp chi tiết được gom thành 20 siêu lớp, giúp phân tích lỗi kiểu
        ``nhầm trong cùng nhóm'' hay ``nhầm khác nhóm''.
\end{itemize}

\subsection{Nguồn và license}
CIFAR-100 do Krizhevsky (2009) công bố tại \url{https://www.cs.toronto.edu/~kriz/cifar.html}, là tập con có gán nhãn
của bộ ``80 Million Tiny Images''. Tập không đi kèm license chuẩn (kiểu MIT hay CC); tác giả cho phép dùng tự do cho
nghiên cứu và học tập, với yêu cầu trích dẫn báo cáo kỹ thuật gốc~\cite{krizhevsky2009}. Bộ 80 Million Tiny Images
gốc đã bị gỡ năm 2020 do chứa nhãn không phù hợp, nhưng CIFAR-10/100 là tập đã được gán nhãn lại thủ công và vẫn được
cung cấp chính thức.

Do server gốc tải rất chậm từ Việt Nam, nhóm tải dữ liệu từ bản sao chính thức trên Hugging Face
(\code{uoft-cs/cifar100}). Bản này có cùng ảnh, cùng nhãn và cùng thứ tự với bản gốc. Ảnh được giải mã một lần
và cache vào \code{data/cifar100.npz} (\code{data.py}); thư mục \code{data/} không được commit lên GitHub.
Metadata của bản Hugging Face ghi sai tên lớp \code{crab} thành \code{cra}; \code{data.py} sửa lại tên gốc khi nạp.

\subsection{Cách chia split}
\begin{table}[H]
\centering
\caption{Ba tập dữ liệu dùng chung cho cả 8 thí nghiệm (\code{output/eda/summary.json}).}\label{tab:split}
\begin{tabularx}{\linewidth}{@{}lrrlL@{}}
\toprule
Split & Số ảnh & Mỗi lớp & Nguồn & Vai trò \\\midrule
Train      & 45.000 & 450 & Tập train gốc & Cập nhật trọng số \\
Validation &  5.000 &  50 & Tập train gốc & Chọn checkpoint tốt nhất, theo dõi quá trình học \\
Test       & 10.000 & 100 & Tập test gốc  & Chỉ dùng một lần ở cuối mỗi thí nghiệm \\
\bottomrule
\end{tabularx}
\end{table}

Validation được lấy \st{phân tầng}: mỗi lớp chọn ngẫu nhiên 50 ảnh từ tập train gốc với seed 42
(\code{make\_split} trong \code{data.py}). Chỉ số ảnh được lưu cố định trong \code{splits/cifar100\_split.json},
nên cả 8 thí nghiệm dùng \st{đúng cùng một split}. Tập test không được dùng để chọn siêu tham số hay checkpoint.

\subsection{Khám phá dữ liệu (EDA)}
Script \code{eda.py} đếm số ảnh mỗi lớp theo split, vẽ ảnh mẫu, ví dụ augmentation, tính thống kê pixel và
liệt kê 20 siêu lớp; kết quả lưu ở \code{output/eda/}.

\begin{figure}[H]
\centering
\fig{\linewidth}{class_counts.png}
\caption{Số ảnh mỗi lớp theo split: dữ liệu cân bằng tuyệt đối (450 / 50 / 100 ảnh mỗi lớp).}\label{fig:counts}
\end{figure}

\begin{figure}[H]
\centering
\fig{0.82\linewidth}{samples_per_class.png}
\caption{Một ảnh mẫu cho mỗi lớp của CIFAR-100 (độ phân giải gốc $32\times32$).}\label{fig:samples}
\end{figure}

\textbf{Nhận xét.}
\begin{itemize}
  \item Mọi ảnh là ảnh màu RGB kích thước $32\times32\times3$. Không lớp nào thiếu ảnh (Hình~\ref{fig:counts}),
        nên accuracy và macro-F1 phản ánh tương đương nhau; nhóm vẫn báo cáo cả hai.
  \item Trung bình pixel trên tập train (R, G, B) là (0.507, 0.487, 0.441), độ lệch chuẩn là (0.267, 0.256, 0.276).
        Thống kê này gần với thống kê ImageNet (mean (0.485, 0.456, 0.406), std (0.229, 0.224, 0.225)), nên việc dùng
        chuẩn hoá ImageNet cho mọi cấu hình là hợp lý (Mục~\ref{sec:loader}).
  \item Ảnh rất nhỏ, vật thể thường chiếm gần hết khung hình và nhiều chi tiết bị mất. Lỗi dự kiến tập trung ở các
        lớp cùng siêu lớp (Bảng~\ref{tab:coarse}), ví dụ các loài cây, các loại xe hay các loài thú có vú cỡ nhỏ.
\end{itemize}

\begin{xltabular}{\linewidth}{@{}P{5.2cm}L@{}}
\caption{20 siêu lớp của CIFAR-100, mỗi siêu lớp gồm 5 lớp chi tiết.}\label{tab:coarse}\\
\toprule Siêu lớp & Các lớp chi tiết \\\midrule
\endfirsthead
\toprule Siêu lớp & Các lớp chi tiết \\\midrule
\endhead
\bottomrule
\endlastfoot
aquatic mammals & beaver, dolphin, otter, seal, whale \\
fish & aquarium fish, flatfish, ray, shark, trout \\
flowers & orchid, poppy, rose, sunflower, tulip \\
food containers & bottle, bowl, can, cup, plate \\
fruit and vegetables & apple, mushroom, orange, pear, sweet pepper \\
household electrical devices & clock, keyboard, lamp, telephone, television \\
household furniture & bed, chair, couch, table, wardrobe \\
insects & bee, beetle, butterfly, caterpillar, cockroach \\
large carnivores & bear, leopard, lion, tiger, wolf \\
large man-made outdoor things & bridge, castle, house, road, skyscraper \\
large natural outdoor scenes & cloud, forest, mountain, plain, sea \\
large omnivores and herbivores & camel, cattle, chimpanzee, elephant, kangaroo \\
medium mammals & fox, porcupine, possum, raccoon, skunk \\
non-insect invertebrates & crab, lobster, snail, spider, worm \\
people & baby, boy, girl, man, woman \\
reptiles & crocodile, dinosaur, lizard, snake, turtle \\
small mammals & hamster, mouse, rabbit, shrew, squirrel \\
trees & maple tree, oak tree, palm tree, pine tree, willow tree \\
vehicles 1 & bicycle, bus, motorcycle, pickup truck, train \\
vehicles 2 & lawn mower, rocket, streetcar, tank, tractor \\
\end{xltabular}

% ======================================================================= 3
\section{DataLoader và augmentation}\label{sec:loader}
\begin{table}[H]
\centering
\caption{Pipeline tiền xử lý, giống hệt nhau cho cả 8 cấu hình (\code{data.py}).}\label{tab:pipeline}
\begin{tabularx}{\linewidth}{@{}P{3.6cm}LL@{}}
\toprule
Bước & Train & Validation / Test \\\midrule
1. Augmentation ở $32\times32$ & \code{RandomCrop(32, padding=4)} + \code{RandomHorizontalFlip} & Không \\
2. Phóng ảnh & \multicolumn{2}{l}{Nội suy bicubic lên $224\times224$ (chạy trên GPU, xem Mục~\ref{sec:env})} \\
3. Chuẩn hoá & \multicolumn{2}{l}{Chuẩn hoá theo mean/std của ImageNet} \\
DataLoader & batch 128, \code{shuffle}, \code{drop\_last} & batch 128, không shuffle \\
\bottomrule
\end{tabularx}
\end{table}

\begin{figure}[H]
\centering
\fig{\linewidth}{augmentation_examples.png}
\caption{Một ảnh gốc và các phiên bản sau augmentation (đã phóng lên $224\times224$).}\label{fig:aug}
\end{figure}

\textbf{Lý do thiết kế.}
\begin{itemize}
  \item \st{Phóng lên $224\times224$:} trọng số pretrained của cả hai mô hình được học ở $224\times224$. Riêng DeiT-Tiny
        dùng patch $16\times16$; nếu giữ ảnh $32\times32$ thì chỉ còn 4 token, không đủ để attention hoạt động có ý nghĩa.
  \item \st{Cùng kích thước, cùng augmentation, cùng chuẩn hoá cho cả 8 cấu hình} (kể cả from-scratch): khác biệt kết quả
        chỉ đến từ ba yếu tố cần so sánh là kiến trúc, cách khởi tạo và mức đóng băng.
  \item \st{Augmentation nhẹ:} crop + flip là chuẩn phổ biến cho CIFAR. Nhóm cố ý không dùng augmentation mạnh
        (RandAugment, Mixup, \ldots) để không ưu ái riêng Transformer from-scratch, vốn hưởng lợi nhiều từ augmentation mạnh.
        Đây là hạn chế có chủ đích và sẽ được thảo luận ở A1.2.
\end{itemize}

% ======================================================================= 4
\section{Mô hình sử dụng}
\begin{table}[H]
\centering
\caption{Hai mô hình đại diện cho hai họ kiến trúc (\code{config.py}, \code{models.py}).}\label{tab:models}
\begin{tabularx}{\linewidth}{@{}P{2.5cm}P{4.6cm}P{2.2cm}L@{}}
\toprule
Họ & Mô hình (tên trong \code{timm}) & Tham số (head 100 lớp) & Lý do chọn \\\midrule
CNN & ResNet-18~\cite{he2016} \newline \code{resnet18.a1\_in1k} & 11.227.812 & CNN chuẩn, nhỏ, có kết nối tắt (residual); baseline phổ biến nhất \\
Transformer & DeiT-Tiny/16~\cite{touvron2021} \newline \code{deit\_tiny\_patch16\_224}\newline\code{.fb\_in1k} & 5.543.716 & Một trong các ViT nhỏ nhất có trọng số ImageNet-1k tốt, cùng cỡ với ResNet-18, chạy được trên GPU phổ thông \\
\bottomrule
\end{tabularx}
\end{table}

Hai mô hình có cùng nguồn pretrained (chỉ ImageNet-1k, không dùng ImageNet-21k), lấy qua thư viện
\code{timm}~\cite{timm}, và cùng bậc kích thước, để so sánh không bị lệch vì một bên được học trước trên nhiều dữ liệu hơn.
Lớp phân loại cuối của cả hai được thay bằng lớp mới 100 đầu ra, khởi tạo ngẫu nhiên.

% ======================================================================= 5
\section{Protocol thí nghiệm}
\subsection{Ma trận 8 thí nghiệm}
\begin{table}[H]
\centering
\caption{8 cấu hình và số tham số trainable, đo bằng \code{smoke\_test.py} (\code{output/protocol\_table.json}).}\label{tab:matrix}
\begin{tabularx}{\linewidth}{@{}lllLrr@{}}
\toprule
Cấu hình & Kiến trúc & Khởi tạo & Chế độ & Trainable & Tỉ lệ \\\midrule
\code{resnet18\_scratch}  & ResNet-18 & Ngẫu nhiên  & Học từ đầu         & 11.227.812 & 100,0\% \\
\code{resnet18\_head}     & ResNet-18 & ImageNet-1k & Chỉ train head     &     51.300 &   0,5\% \\
\code{resnet18\_partial}  & ResNet-18 & ImageNet-1k & Đóng băng một phần & 10.544.740 &  93,9\% \\
\code{resnet18\_full}     & ResNet-18 & ImageNet-1k & Fine-tune toàn bộ  & 11.227.812 & 100,0\% \\\midrule
\code{deit\_tiny\_scratch} & DeiT-Tiny & Ngẫu nhiên  & Học từ đầu         &  5.543.716 & 100,0\% \\
\code{deit\_tiny\_head}    & DeiT-Tiny & ImageNet-1k & Chỉ train head     &     19.300 &   0,4\% \\
\code{deit\_tiny\_partial} & DeiT-Tiny & ImageNet-1k & Đóng băng một phần &  2.688.868 &  48,5\% \\
\code{deit\_tiny\_full}    & DeiT-Tiny & ImageNet-1k & Fine-tune toàn bộ  &  5.543.716 & 100,0\% \\
\bottomrule
\end{tabularx}
\end{table}

\textbf{Định nghĩa các mức đóng băng} (áp dụng trên bản pretrained, \code{models.py}):
\begin{itemize}
  \item \st{Chỉ train head} (linear probe): đóng băng toàn bộ backbone, chỉ học lớp phân loại cuối.
  \item \st{Đóng băng một phần:} khoảng nửa đầu của mạng. ResNet-18 đóng băng \code{conv1}, \code{bn1}, \code{layer1},
        \code{layer2} và train \code{layer3}, \code{layer4}, \code{fc}. DeiT-Tiny đóng băng \code{patch\_embed},
        positional embedding, token CLS và 6/12 block đầu; train 6 block cuối, LayerNorm cuối và head.
  \item \st{Fine-tune toàn bộ:} mọi tham số đều được cập nhật.
\end{itemize}

Hai mô hình được đóng băng cùng ``một nửa độ sâu'' (2/4 stage của ResNet-18, 6/12 block của DeiT-Tiny), nhưng tỉ lệ
tham số trainable khác nhau: 93,9\% so với 48,5\%. Lý do là ResNet tập trung phần lớn tham số ở các stage sâu
(\code{layer3}, \code{layer4}), còn các block của DeiT có kích thước bằng nhau. Nhóm sẽ báo cáo cả hai góc nhìn
(độ sâu và số tham số) khi thảo luận ở A1.2. Với ResNet-18, các lớp BatchNorm thuộc phần bị đóng băng được giữ ở chế
độ \code{eval} trong lúc train (\code{set\_train\_mode}), để thống kê pretrained không bị thay đổi.

\subsection{Huấn luyện}
\begin{table}[H]
\centering
\caption{Siêu tham số huấn luyện, dùng đúng như kế hoạch A1.1 (\code{config.TRAIN\_HPARAMS}).}\label{tab:hparams}
\begin{tabularx}{\linewidth}{@{}P{3.6cm}LL@{}}
\toprule
Thiết lập & From-scratch & Pretrained (3 chế độ) \\\midrule
Optimizer & \multicolumn{2}{l}{AdamW, weight decay 0.05} \\
Learning rate & $10^{-3}$ & $10^{-4}$ \\
Lịch learning rate & \multicolumn{2}{l}{Warmup tuyến tính 1 epoch, sau đó cosine giảm về 0} \\
Số epoch & 30 & 10 \\
Loss & \multicolumn{2}{l}{Cross-entropy, label smoothing 0.1} \\
Batch size / seed & \multicolumn{2}{l}{128 / 42} \\
Chọn checkpoint & \multicolumn{2}{P{11.5cm}}{Epoch có validation accuracy cao nhất; sau đó đánh giá trên test đúng một lần} \\
\bottomrule
\end{tabularx}
\end{table}

Vòng lặp huấn luyện được tự viết bằng PyTorch trong \code{train.py}: xoá gradient $\rightarrow$ forward $\rightarrow$
tính loss $\rightarrow$ backward $\rightarrow$ \code{optimizer.step()}. Thiết bị được chọn tự động theo thứ tự CUDA, MPS, CPU.
From-scratch được học nhiều epoch hơn và với learning rate lớn hơn, vì mô hình phải học mọi thứ từ đầu.
Để so sánh công bằng trong từng nhóm, các cấu hình cùng nhóm dùng chung siêu tham số.

\subsection{Metric và câu hỏi nghiên cứu}
\begin{itemize}
  \item \st{Chất lượng:} top-1 accuracy và macro-F1 (\code{sklearn.metrics.f1\_score}) trên validation và test;
        đường cong loss/accuracy theo epoch (\code{train.py} đã ghi đủ các giá trị này vào log JSON của mỗi lần chạy).
  \item \st{Chi phí:} thời gian huấn luyện, số tham số trainable, bộ nhớ GPU đỉnh (\code{torch.cuda.max\_memory\_allocated}).
  \item \st{Phân tích lỗi:} confusion matrix ở mức 20 siêu lớp (dễ đọc hơn ma trận $100\times100$), các cặp lớp hay nhầm
        nhất và ảnh mẫu bị đoán sai của cấu hình tốt nhất mỗi họ. Dự đoán trên test được lưu sẵn
        (\code{*\_test\_preds.json}) để phục vụ phân tích này.
\end{itemize}

\textbf{Câu hỏi cần trả lời và giả thuyết ban đầu} (đặt ra ở A1.1, trước khi có số liệu; được kiểm chứng ở
Mục~\ref{sec:discussion}):
\begin{enumerate}
  \item \st{ViT có cần pretrain không?} Dự đoán DeiT-Tiny from-scratch kém hơn rõ rệt so với ResNet-18 from-scratch,
        vì Transformer thiếu ``thiên kiến cục bộ'' của tích chập và CIFAR-100 chỉ có 45.000 ảnh train.
  \item \st{Khi nào đóng băng có lợi?} Dự đoán đóng băng một phần đạt gần bằng fine-tune toàn bộ với ít tham số trainable
        hơn. Chỉ train head có thể kém hơn, vì ảnh $32\times32$ phóng lên khác khá xa ảnh ImageNet.
  \item \st{Chi phí so với chất lượng:} so sánh accuracy trên mỗi đơn vị thời gian huấn luyện giữa các cấu hình.
\end{enumerate}

% ======================================================================= 6
\section{Kiểm tra pipeline ở mốc A1.1}
Script \code{smoke\_test.py} tạo đủ 8 cấu hình, đếm tham số và chạy một bước forward + backward trên một batch 8 ảnh
ngẫu nhiên $224\times224$. Script xác nhận các tham số bị đóng băng không nhận gradient, còn các tham số trainable thì có
gradient. Cả 8 cấu hình đều đạt (cột ``Kiểm tra'' = OK trong \code{output/protocol\_table.md}; số tham số ở
Bảng~\ref{tab:matrix}).

Ngoài ra, nhóm chạy thử toàn bộ vòng lặp huấn luyện -- đánh giá trên tập con nhỏ (256 ảnh mỗi tập, 1 epoch, thiết bị
\code{mps}) để chắc chắn pipeline chạy thông: \code{deit\_tiny\_partial} (khoảng 6 giây) và \code{resnet18\_head}
(khoảng 3 giây). Huấn luyện, đánh giá validation, chọn checkpoint và đánh giá test đều chạy không lỗi; log nằm ở
\code{output/smoke/}. Accuracy của các lần chạy này chỉ ở mức đoán ngẫu nhiên (khoảng 1--3\%) vì chỉ có 1 epoch trên
256 ảnh: đây là kiểm tra code, \st{không phải kết quả thí nghiệm}.
% ======================================================================= 7
\section{Môi trường chạy A1.2}\label{sec:env}
Cả 8 thí nghiệm được chạy trên Google Colab với GPU \val{GpuName}, bằng notebook \code{colab\_a1\_2.ipynb}
(gọi \code{run\_all.py}, lần lượt chạy \code{train.py} cho từng cấu hình). Thứ tự chạy: 6 cấu hình pretrained trước,
2 cấu hình from-scratch sau. Máy Mac M1 8\,GB của nhóm không đủ cho protocol $224\times224$: đo thử, một epoch
ResNet-18 from-scratch mất khoảng 17 phút, còn DeiT-Tiny from-scratch vượt quá bộ nhớ và chậm hơn khoảng 100 lần.

So với mô tả ở A1.1, nhóm \st{giữ nguyên protocol} (split, augmentation, kích thước ảnh, mô hình, mức đóng băng,
siêu tham số ở Bảng~\ref{tab:hparams}) và chỉ thay đổi ba chi tiết cài đặt để chạy được trên Colab:
\begin{itemize}
  \item \st{Phóng ảnh trên GPU.} DataLoader chỉ làm crop + lật ở $32\times32$; bước phóng bicubic lên $224\times224$ và
        chuẩn hoá (\code{to\_model\_input} trong \code{data.py}) chạy trên GPU, vì 2 lõi CPU của Colab không phóng ảnh kịp
        cho GPU. Kết quả lệch so với cách phóng bằng PIL trung bình khoảng 1 mức xám trên 255 (PIL làm tròn về số nguyên
        8 bit và dùng hệ số bicubic hơi khác), nên không ảnh hưởng tới so sánh.
  \item \st{Mixed precision (AMP, fp16)}~\cite{micikevicius2018} cho forward/backward, kèm \code{GradScaler} để tránh
        gradient quá nhỏ bị làm tròn về 0. Áp dụng như nhau cho cả 8 cấu hình.
  \item \st{Checkpoint sau mỗi epoch.} Phiên Colab miễn phí có thể bị ngắt; chạy lại cùng lệnh sẽ tiếp tục từ epoch cuối
        (trạng thái mô hình, optimizer, lịch learning rate, bộ sinh số ngẫu nhiên đều được lưu). Nhóm đã kiểm tra: một lần
        chạy bị dừng giữa chừng rồi chạy tiếp cho kết quả trùng khớp từng epoch với lần chạy liền một mạch.
\end{itemize}

\textbf{Cách đo chi phí.} Thời gian huấn luyện là tổng thời gian các epoch train (không tính đánh giá validation),
đo với \code{torch.cuda.synchronize()}; bộ nhớ là \code{torch.cuda.max\_memory\_allocated()}; tốc độ dự đoán là số ảnh
test xử lý mỗi giây ở batch 128 (đã gồm bước phóng ảnh, bỏ batch đầu tiên để GPU ``khởi động'').

% ======================================================================= 8
\section{Kết quả}\label{sec:results}
\subsection{Bảng tổng hợp}
Mỗi cấu hình được đánh giá trên test đúng một lần, bằng checkpoint có validation accuracy cao nhất.

\begin{table}[H]
\centering\small
\caption{Chất lượng của 8 cấu hình (\%). ``Epoch tốt nhất'': epoch có validation accuracy cao nhất / tổng số epoch;
``Train acc'': accuracy trên tập train ở epoch cuối (có augmentation).}\label{tab:quality}
\begin{tabularx}{\linewidth}{@{}lLcrrrr@{}}
\toprule
Cấu hình & Chế độ & \makecell{Epoch\\tốt nhất} & \makecell{Train\\acc} & \makecell{Val\\acc} & \makecell{Test\\acc} & \makecell{Test\\macro-F1} \\\midrule
\gen{table_quality.tex}{7}
\bottomrule
\end{tabularx}
\end{table}

\begin{table}[H]
\centering\small
\caption{Chi phí của 8 cấu hình trên GPU \val{GpuName}.}\label{tab:cost}
\begin{tabularx}{\linewidth}{@{}lLrrrr@{}}
\toprule
Cấu hình & Tham số trainable & \makecell{Train\\(phút)} & \makecell{Giây /\\epoch} & \makecell{Bộ nhớ\\GPU (GB)} & \makecell{Dự đoán\\(ảnh/s)} \\\midrule
\gen{table_cost.tex}{6}
\bottomrule
\end{tabularx}
\end{table}

\begin{figure}[H]
\centering
\fig{0.92\linewidth}{test_acc_by_mode.png}
\caption{Test accuracy theo kiến trúc và chế độ huấn luyện.}\label{fig:acc}
\end{figure}

\textbf{Nhận xét.}
\begin{itemize}
  \item \st{Thứ hạng trong mỗi họ.} Với cả hai họ, fine-tune toàn bộ đứng đầu, đóng băng một phần đứng thứ hai.
        Hai họ khác nhau ở hai chế độ còn lại: với ResNet-18, học từ đầu (\val{ResNetScratchTest}\%) tốt hơn hẳn chỉ
        train head (\val{ResNetHeadTest}\%); với DeiT-Tiny thì ngược lại, chỉ train head (\val{DeiTHeadTest}\%) tốt hơn
        học từ đầu (\val{DeiTScratchTest}\%).
  \item \st{Khoảng cách giữa hai họ.} Khi học từ đầu, ResNet-18 hơn DeiT-Tiny 19,0 điểm. Khi dùng trọng số pretrained,
        DeiT-Tiny hơn ResNet-18 ở mọi mức đóng băng: +14,1 điểm (chỉ head), +7,6 (một phần), +10,3 (toàn bộ), dù
        DeiT-Tiny chỉ có khoảng một nửa số tham số (5,5 triệu so với 11,2 triệu).
  \item \st{Macro-F1 bám sát accuracy} (chênh lệch tối đa 1,8 điểm, ở \code{resnet18\_head}) vì test có đúng 100 ảnh mỗi
        lớp. Test accuracy cũng lệch validation accuracy không quá 1,1 điểm ở mọi cấu hình, nên việc chọn checkpoint theo
        validation không làm kết quả test bị lạc quan.
\end{itemize}

\subsection{Quá trình học}
\begin{figure}[H]
\centering
\fig{0.95\linewidth}{learning_curves.png}
\caption{Validation accuracy (hàng trên) và train accuracy (hàng dưới) theo epoch; số ở cuối mỗi đường là giá trị epoch cuối.}\label{fig:curves}
\end{figure}

\textbf{Nhận xét.}
\begin{itemize}
  \item \st{Tốc độ hội tụ.} Các cấu hình pretrained đạt gần mức cuối chỉ sau 3--5 epoch (DeiT-Tiny đóng băng một phần
        hoặc toàn bộ đã có khoảng 76\% validation accuracy sau epoch 2). Hai cấu hình học từ đầu cần khoảng 20 epoch mới
        chững lại; epoch tốt nhất là \val{ResNetScratchBestEpoch} (ResNet-18) và \val{DeiTScratchBestEpoch} (DeiT-Tiny).
  \item \st{Overfitting khi học từ đầu.} Train accuracy ở epoch cuối gần 100\% (ResNet-18 \val{ResNetScratchTrainAcc}\%,
        DeiT-Tiny \val{DeiTScratchTrainAcc}\%) trong khi validation dừng ở khoảng 68\% và 48\%. Với DeiT-Tiny, validation
        gần như đứng yên từ epoch 16 trong khi train vẫn tăng thêm khoảng 40 điểm: mô hình học thuộc tập train thay vì học
        đặc trưng tổng quát.
  \item \st{Pretrained ít overfit hơn.} ResNet-18 đóng băng một phần hoặc toàn bộ có train accuracy chỉ cao hơn validation
        khoảng 2 điểm. Train accuracy của chúng chỉ khoảng 77\%, tức mô hình chưa khớp hết tập train: ResNet-18 pretrained có
        thể còn tăng nếu train lâu hơn hoặc dùng learning rate lớn hơn $10^{-4}$. DeiT-Tiny fine-tune toàn bộ có khoảng cách lớn hơn (\val{DeiTFullTrainAcc}\% so với
        \val{DeiTFullVal}\%) nhưng validation vẫn tăng nhẹ đến epoch cuối (epoch tốt nhất là \val{DeiTFullBestEpoch}).
  \item \st{Chỉ train head bị chặn trên.} Train và validation accuracy gần bằng nhau và đi ngang sau khoảng epoch 6: một lớp
        tuyến tính trên đặc trưng cố định không đủ biểu diễn 100 lớp CIFAR-100 (underfitting chứ không phải overfitting).
\end{itemize}

\subsection{Chất lượng so với chi phí}
\begin{figure}[H]
\centering
\fig{0.95\linewidth}{accuracy_vs_cost.png}
\caption{Test accuracy theo thời gian huấn luyện (trái) và theo số tham số trainable (phải, thang log).
Màu là chế độ huấn luyện, hình tròn là ResNet-18, hình vuông là DeiT-Tiny.}\label{fig:cost}
\end{figure}

\textbf{Nhận xét.}
\begin{itemize}
  \item \st{Thời gian.} Hai cấu hình học từ đầu tốn nhiều thời gian nhất (33--37 phút, 30 epoch) nhưng lại thuộc nhóm
        kém nhất. \code{deit\_tiny\_full} vừa chính xác nhất vừa nhanh hơn \code{resnet18\_full}
        (\val{DeiTFullMin} so với \val{ResNetFullMin} phút). \code{deit\_tiny\_partial} đạt \val{DeiTPartialTest}\% chỉ trong
        \val{DeiTPartialMin} phút, là lựa chọn tốt nếu thời gian bị giới hạn.
  \item \st{Số tham số trainable không phản ánh hết chi phí.} \code{resnet18\_partial} vẫn train 93,9\% tham số nhưng nhanh
        gấp đôi \code{resnet18\_full} và tốn chưa đến một nửa bộ nhớ, vì các tầng đầu bị đóng băng là các tầng có feature map
        lớn nhất (từ $112\times112$ xuống $28\times28$): bỏ được lan truyền ngược qua chúng mới là phần tiết kiệm chính.
  \item \st{Bộ nhớ GPU} tăng theo phần mô hình phải lan truyền ngược: chỉ train head dùng 0,25--0,60\,GB, fine-tune toàn bộ
        dùng 1,7--2,2\,GB. Cả 8 cấu hình đều vừa GPU T4 16\,GB với batch 128.
  \item \st{Tốc độ dự đoán} gần như bằng nhau (khoảng 1.900 ảnh/s) ở mọi cấu hình, vì chế độ huấn luyện không đổi kiến trúc
        lúc dự đoán, và ở batch 128 bước nạp dữ liệu và phóng ảnh chiếm phần lớn thời gian. Hai cấu hình chỉ train head có
        cùng 23 giây mỗi epoch vì cùng lý do.
\end{itemize}

% ======================================================================= 9
\section{Phân tích lỗi}\label{sec:errors}
Phần này so sánh cấu hình tốt nhất của mỗi họ, chọn theo validation accuracy: \val{BestResNetName} (ResNet-18) và
\val{BestDeiTName} (DeiT-Tiny).

\subsection{Nhầm lẫn theo siêu lớp}
Confusion matrix $100\times100$ quá lớn để đọc, nên nhóm gom dự đoán về 20 siêu lớp (Bảng~\ref{tab:coarse}).
Accuracy ở mức siêu lớp: ResNet-18 \val{BestResNetCoarseAcc}\%, DeiT-Tiny \val{BestDeiTCoarseAcc}\%.
Trong số các ảnh bị đoán sai lớp chi tiết, tỉ lệ ảnh bị đoán nhầm sang một lớp \emph{cùng siêu lớp} là
\val{BestResNetSameSuper}\% với ResNet-18 và \val{BestDeiTSameSuper}\% với DeiT-Tiny (nếu nhầm ngẫu nhiên, tỉ lệ này
chỉ khoảng $4/99 \approx 4\%$).

\begin{figure}[H]
\centering
\fig{0.78\linewidth}{confusion_coarse_resnet18.png}
\caption{Confusion matrix theo 20 siêu lớp của ResNet-18 tốt nhất (mỗi hàng chuẩn hoá về 100\%).}\label{fig:cm-resnet}
\end{figure}

\begin{figure}[H]
\centering
\fig{0.78\linewidth}{confusion_coarse_deit_tiny.png}
\caption{Confusion matrix theo 20 siêu lớp của DeiT-Tiny tốt nhất.}\label{fig:cm-deit}
\end{figure}

\begin{figure}[H]
\centering
\fig{0.85\linewidth}{superclass_gap.png}
\caption{Chênh lệch accuracy theo siêu lớp giữa DeiT-Tiny và ResNet-18 tốt nhất (dương: DeiT-Tiny tốt hơn).}\label{fig:gap}
\end{figure}

\subsection{Các cặp lớp hay nhầm nhất}
\begin{table}[H]
\centering\small
\caption{8 cặp (lớp thật $\rightarrow$ lớp dự đoán) bị nhầm nhiều nhất trên 10.000 ảnh test (mỗi lớp có 100 ảnh).}\label{tab:pairs}
\begin{minipage}[t]{0.48\linewidth}\centering
ResNet-18 (\val{BestResNetName})\\[2pt]
\begin{tabular}{@{}llr@{}}
\toprule Thật & Dự đoán & Số ảnh \\\midrule
\gen{confused_pairs_resnet18.tex}{3}
\bottomrule
\end{tabular}
\end{minipage}\hfill
\begin{minipage}[t]{0.48\linewidth}\centering
DeiT-Tiny (\val{BestDeiTName})\\[2pt]
\begin{tabular}{@{}llr@{}}
\toprule Thật & Dự đoán & Số ảnh \\\midrule
\gen{confused_pairs_deit_tiny.tex}{3}
\bottomrule
\end{tabular}
\end{minipage}
\end{table}

Năm lớp khó nhất (accuracy theo lớp): ResNet-18: \val{BestResNetWorst}; DeiT-Tiny: \val{BestDeiTWorst}.
Năm lớp dễ nhất: ResNet-18: \val{BestResNetTop}; DeiT-Tiny: \val{BestDeiTTop}.

\subsection{Hai mô hình có sai cùng chỗ không?}
Trên 10.000 ảnh test: cả hai cùng đúng \val{AgreeBothRight}\%, chỉ ResNet-18 đúng \val{AgreeOnlyCnn}\%,
chỉ DeiT-Tiny đúng \val{AgreeOnlyVit}\%, cả hai cùng sai \val{AgreeBothWrong}\%. Nếu luôn chọn được mô hình đúng
trong hai, accuracy sẽ là \val{OracleAcc}\% (giới hạn trên của việc kết hợp hai mô hình).

\begin{figure}[H]
\centering
\fig{0.95\linewidth}{error_examples.png}
\caption{Ảnh test bị đoán sai, chọn ngẫu nhiên (seed 42) theo ba nhóm. ``R'' là dự đoán của ResNet-18, ``D'' là của DeiT-Tiny.}\label{fig:errors}
\end{figure}

\textbf{Nhận xét.}
\begin{itemize}
  \item \st{Lỗi tập trung trong cùng siêu lớp.} Gần một nửa số ảnh bị đoán sai bị nhầm sang lớp cùng siêu lớp, gấp hơn
        10 lần mức ngẫu nhiên. Các cặp hay nhầm nhất ở cả hai mô hình gần như trùng nhau và đều là những lớp giống nhau
        về nghĩa: \emph{girl / woman}, \emph{boy / man}, \emph{maple tree / oak tree}, \emph{seal / otter}. Ở độ phân giải
        $32\times32$, phân biệt tuổi và giới tính của người hay loài cây chỉ dựa vào vài điểm ảnh, nên đây là lỗi khó
        tránh, không phải lỗi ngẫu nhiên.
  \item \st{Lớp khó.} Nhóm \emph{people} (boy, girl, man, woman) nằm trong số lớp khó nhất của cả hai mô hình. Lớp dễ là
        các vật có hình dạng hoặc màu sắc đặc trưng (motorcycle, road, tractor, wardrobe).
  \item \st{DeiT-Tiny tốt hơn ở cả 20 siêu lớp} (Hình~\ref{fig:gap}). Chênh lệch lớn nhất ở các nhóm động vật
        (aquatic mammals +16,2, medium mammals +12,4, small mammals +11,8, fish +10,8 điểm), nhỏ nhất ở đồ vật trong nhà
        (food containers +1,8, household furniture và electrical devices +2,8). Nhóm chưa kiểm chứng nguyên nhân; một giả
        thuyết là các lớp động vật cần phân biệt kết cấu lông, da và bối cảnh trên toàn ảnh, phù hợp với attention toàn cục,
        còn đồ vật có hình dạng rõ thì cả hai mô hình đều đã làm tốt.
  \item \st{Hai mô hình bổ sung cho nhau một phần.} Có \val{AgreeOnlyCnn}\% ảnh chỉ ResNet-18 đúng, nên nếu kết hợp lý
        tưởng hai mô hình thì accuracy có thể lên \val{OracleAcc}\% (so với \val{DeiTFullTest}\% của DeiT-Tiny). Đây chỉ là
        giới hạn trên; nhóm không thử ensemble trong bài này.
  \item \st{Ảnh cả hai cùng sai} (\val{AgreeBothWrong}\%, hàng dưới của Hình~\ref{fig:errors}) thường là ảnh mà vật chính
        bị che, chụp cận một phần, góc nhìn lạ hoặc lẫn vào nền (ví dụ con tôm hùm trong tay người, cận cảnh một góc tủ
        quần áo). Một số ảnh khó cả với người nhìn ở $32\times32$; cải thiện nhóm này cần dữ liệu độ phân giải cao hơn
        hơn là đổi kiến trúc.
\end{itemize}

% ======================================================================= 10
\section{Thảo luận}\label{sec:discussion}
\subsection{ViT có cần pretrain không?}
\st{Giả thuyết 1 được xác nhận.} Học từ đầu với cùng 30 epoch và cùng augmentation, DeiT-Tiny chỉ đạt
\val{DeiTScratchTest}\%, kém ResNet-18 (\val{ResNetScratchTest}\%) 19,0 điểm, và overfit nặng hơn (Hình~\ref{fig:curves}).
Pretrain thay đổi hoàn toàn thứ tự: DeiT-Tiny tăng 36,5 điểm khi dùng trọng số pretrained (từ \val{DeiTScratchTest}\% lên
\val{DeiTFullTest}\%), trong khi ResNet-18 chỉ tăng 7,3 điểm (từ \val{ResNetScratchTest}\% lên \val{ResNetFullTest}\%).
Thậm chí DeiT-Tiny chỉ train head (\val{DeiTHeadTest}\%, 19.300 tham số) vẫn tốt hơn DeiT-Tiny học từ đầu 5,5 triệu tham số.
Kết quả này khớp với lập luận: tích chập có sẵn giả định ``điểm ảnh gần nhau liên quan với nhau'', nên học được từ
45.000 ảnh; Transformer phải tự học quan hệ không gian từ dữ liệu, và CIFAR-100 không đủ lớn để làm việc đó trong 30
epoch với augmentation nhẹ. Như đã ghi ở phần Hạn chế, đây là kết luận cho thiết lập của bài, không phải giới hạn tuyệt
đối của kiến trúc.

\subsection{Khi nào đóng băng có lợi?}
\st{Giả thuyết 2 đúng với ResNet-18, chỉ đúng một phần với DeiT-Tiny.}
\begin{itemize}
  \item \st{ResNet-18:} đóng băng một phần (\val{ResNetPartialTest}\%) chỉ kém fine-tune toàn bộ (\val{ResNetFullTest}\%)
        1,0 điểm, trong khi thời gian train giảm một nửa (\val{ResNetPartialMin} so với \val{ResNetFullMin} phút) và bộ nhớ
        giảm từ \val{ResNetFullMem} xuống \val{ResNetPartialMem}\,GB. Đây là lựa chọn hợp lý khi tài nguyên hạn chế.
        Tuy nhiên tiết kiệm đến từ việc bỏ lan truyền ngược qua các tầng đầu, không phải từ số tham số (vẫn train 93,9\%).
  \item \st{DeiT-Tiny:} đóng băng 6 khối đầu (\val{DeiTPartialTest}\%) kém fine-tune toàn bộ (\val{DeiTFullTest}\%) 3,7
        điểm, đổi lại train nhanh hơn khoảng 35\% và tốn khoảng một nửa bộ nhớ, với 48,5\% tham số trainable. Các khối đầu
        của DeiT-Tiny dường như vẫn cần điều chỉnh cho ảnh CIFAR phóng to, khác với các tầng đầu của CNN vốn học đặc trưng
        cấp thấp (cạnh, màu) dùng chung được.
  \item \st{Chỉ train head kém xa ở cả hai họ} (thấp hơn fine-tune toàn bộ 23,7 điểm với ResNet-18 và 19,9 điểm với
        DeiT-Tiny), đúng như dự đoán: đặc trưng ImageNet cố định không khớp đủ với ảnh $32\times32$ phóng lên. Chế độ này
        chỉ đáng dùng khi cần huấn luyện cực nhanh hoặc rất ít bộ nhớ.
\end{itemize}

\subsection{Chi phí so với chất lượng}
\st{Giả thuyết 3 (so sánh accuracy trên mỗi đơn vị thời gian):} trên GPU T4, fine-tune từ trọng số pretrained tốt hơn học
từ đầu ở cả hai trục: \code{deit\_tiny\_full} đạt accuracy cao nhất (\val{DeiTFullTest}\%) với khoảng một phần ba thời gian
của các cấu hình học từ đầu. Thứ tự đề xuất theo ngân sách:
\begin{itemize}
  \item Khoảng 4 phút: \code{deit\_tiny\_head} (\val{DeiTHeadTest}\%), không nên dùng \code{resnet18\_head}
        (\val{ResNetHeadTest}\%) với cùng chi phí.
  \item Khoảng 6--7 phút: \code{deit\_tiny\_partial} (\val{DeiTPartialTest}\%) hoặc \code{resnet18\_partial}
        (\val{ResNetPartialTest}\%).
  \item Khoảng 11 phút: \code{deit\_tiny\_full} (\val{DeiTFullTest}\%), cấu hình tốt nhất của bài.
\end{itemize}
Học từ đầu chỉ có ý nghĩa khi không được dùng trọng số pretrained; khi đó nên chọn CNN (ResNet-18) thay vì ViT.

\subsection{Hạn chế}
\begin{itemize}
  \item \st{Một seed cho mỗi cấu hình.} Do giới hạn GPU, mỗi cấu hình chỉ chạy một lần; chênh lệch nhỏ hơn khoảng
        1 điểm phần trăm giữa hai cấu hình nên được xem là chưa chắc chắn.
  \item \st{Siêu tham số không được tinh chỉnh riêng.} Mỗi nhóm cấu hình dùng chung một bộ siêu tham số
        (Bảng~\ref{tab:hparams}); mỗi cấu hình có thể đạt cao hơn nếu được tinh chỉnh riêng.
  \item \st{Augmentation nhẹ, ít epoch cho from-scratch.} Kết quả from-scratch (đặc biệt của DeiT-Tiny) phản ánh đúng
        thiết lập này, không phải giới hạn của kiến trúc: DeiT gốc học từ đầu trên ImageNet với augmentation mạnh và
        300 epoch~\cite{touvron2021}.
  \item \st{Ảnh $32\times32$ phóng lên $224\times224$} không thêm thông tin mới; chi phí tính toán tăng khoảng 49 lần so
        với chạy ở $32\times32$. Nhóm chấp nhận để dùng được trọng số pretrained đúng độ phân giải gốc.
\end{itemize}

% ======================================================================= 11
\section{Tái lập kết quả}
\begin{itemize}
  \item \st{Trên Google Colab:} mở \code{Assignment\_1/colab\_a1\_2.ipynb}, chọn GPU T4 và \emph{Run all}. Log được ghi
        vào Google Drive; nếu phiên bị ngắt, chỉ cần \emph{Run all} lại.
  \item \st{Trên máy có GPU:} \code{pip install -r requirements.txt}, sau đó \code{python run\_all.py} (8 thí nghiệm,
        log ở \code{output/runs/}) và \code{python analyze.py} (bảng, hình và số liệu của báo cáo này).
  \item Split cố định ở \code{splits/cifar100\_split.json}, seed 42 cho mọi thí nghiệm. Log JSON của từng cấu hình
        (lịch sử theo epoch, chi phí, kết quả test) và dự đoán trên test được commit trong \code{output/runs/};
        trọng số (\code{*.pth}) không commit vì dung lượng lớn.
\end{itemize}

% ======================================================================= 12
\section{Sử dụng công cụ AI}
Theo quy định liêm chính học thuật của môn học, nhóm ghi rõ: trợ lý AI (Claude) được dùng để hỗ trợ viết khung mã nguồn
(huấn luyện, phân tích, vẽ biểu đồ), soạn bố cục báo cáo và rà soát lỗi. Mọi số liệu trong báo cáo được sinh tự động từ
log huấn luyện thật bằng \code{analyze.py}, không nhập tay. Nhóm đã đọc, chạy kiểm tra, chỉnh sửa và chịu trách nhiệm
về toàn bộ nội dung; mọi lựa chọn dataset, mô hình, protocol và các nhận xét đều được nhóm hiểu và có thể giải thích.

% ======================================================================= Tài liệu tham khảo
\begin{thebibliography}{9}
\bibitem{krizhevsky2009} A. Krizhevsky. \emph{Learning Multiple Layers of Features from Tiny Images}. Technical report, University of Toronto, 2009.
\bibitem{he2016} K. He, X. Zhang, S. Ren, J. Sun. Deep Residual Learning for Image Recognition. \emph{CVPR}, 2016.
\bibitem{dosovitskiy2021} A. Dosovitskiy et al. An Image is Worth 16x16 Words: Transformers for Image Recognition at Scale. \emph{ICLR}, 2021.
\bibitem{touvron2021} H. Touvron et al. Training data-efficient image transformers \& distillation through attention. \emph{ICML}, 2021.
\bibitem{timm} R. Wightman. PyTorch Image Models (timm). \url{https://github.com/huggingface/pytorch-image-models}.
\bibitem{micikevicius2018} P. Micikevicius et al. Mixed Precision Training. \emph{ICLR}, 2018.
\end{thebibliography}

\end{document}
